\usepackage[inner=3cm, outer=2cm, top=2.5cm, bottom=2.5cm, bindingoffset=1cm]{geometry} % ضبط هامش التجليد

\usepackage{fontspec}
\usepackage{polyglossia}
% \usepackage{luabidi} % It's necessary for RTL with "LuaLaTeX" engine only. If you use "XeLaTeX", comment this line out.
\usepackage{bidi}

\setmainlanguage[numerals=western]{arabic} % For Arabic numbers 1, 2, 3
\setotherlanguage{english}

% \newfontfamily\arabicfont[Script=Arabic, Scale=1.2]{IBM Plex Sans Arabic} % تأكد من وجود هذا الخط في نظامك %(c) Turn this on with LuaLaTeX
\newfontfamily\arabicfont[
    path=./Fonts/, % Assuming your fonts are in a "Fonts" subdirectory
    Script=Arabic,
    Scale=1,
    BoldFont=IBMPlexSansArabic-Bold.ttf, 
    AutoFakeSlant=0.15 %(c) To force italic style for Arabic text, since the font doesn't have a native italic variant.
    ]{IBMPlexSansArabic-Regular.ttf} %(c) Turn this off with XeLaTeX.

\newfontfamily\englishfont[Scale=1]{IBM Plex Sans}

% ---  حزم إضافية واختصارات ---
\usepackage{amsmath, amssymb} % للرموز الرياضية مثل THD
\newcommand{\el}{\textenglish} % اختصار للمصطلحات الإنجليزية

% --- تعريف خط الملاحظات المخصص ---
\newfontfamily\arabicnotefont[
    path=./Fonts/,
    Script=Arabic, 
    Scale=1.2,
    AutoFakeSlant=0.25
    ]{IBMPlexSansArabic-ExtraLight.ttf}
    % --- إنشاء أمر مخصص لتسهيل الكتابة بالملاحظة ---
    \newcommand{\mynote}[1]{{\arabicnotefont\itshape #1}} %(c) \itshape for italic style in notes


\usepackage{fancyhdr}
\pagestyle{fancy} 

% تعديل نظام العناوين برمجياً بشكل مستقل
\renewcommand{\sectionmark}[6]{\markboth{#1}{}} 

% --- 1. تعريف النمط التمهيدي (للفهرس واللاتيني) ---
\fancypagestyle{introstyle}{
    \fancyhf{} 
    \fancyfoot[C]{\fontsize{10}{12}\selectfont \thepage} 
    \renewcommand{\headrulewidth}{0pt}
    \renewcommand{\footrulewidth}{0pt}
}

% --- 2. تعريف النمط الرئيسي (لكامل صلب المشروع وجهين) ---
\fancypagestyle{mainstyle}{
    \fancyhf{} 
    \renewcommand{\headrulewidth}{0pt} 
    \renewcommand{\footrulewidth}{0pt} 

    % RO: الصفحات الفردية (رقم الصفحة | اسم القسم)
    \fancyfoot[RO]{\fontsize{10}{12}\selectfont \thepage\ \ |\ \ \leftmark}

    % LE: الصفحات الزوجية (القسم + رقم القسم | رقم الصفحة) - تمت كتابة العربي مباشرة
    \fancyfoot[LE]{\fontsize{10}{12}\selectfont القسم \thesection\ \ |\ \ \thepage}

    \fancyhfoffset[RO,LE]{0pt}
}


% --- حل مشكلة خروج الكلمات إلى الهامش وتوزيع الأسطر تلقائياً ---
\emergencystretch=3em

\usepackage{pict2e} % For advanced picture drawing (used by fancybox)
\usepackage{fancybox} % For shadow, oval, and double boxes 
\usepackage{xcolor} % For colors

\usepackage{tikz}
\usepackage{varwidth}

\usepackage{graphicx}
\usepackage{float}
\graphicspath{{figures/}}


% 1. تعريف لون مخصص للصناديق الإرشادية لسهولة التعديل لاحقاً
\definecolor{boxbackground}{gray}{0.93} % يعادل gray!15 تقريباً ولكن بنقاء أعلى
\definecolor{boxborder}{gray}{0.2}      % إطار رمادي داكن مريح للعين بدلاً من الأسود الحاد

% تعريف خط مخصص لعنوان الصندوق (وزن شبه عريض SemiBold)
\newfontfamily\boxTitleFont[Script=Arabic]{IBMPlexSansArabic-SemiBold.ttf}
% تعريف خط مخصص لمتن الصندوق (وزن عادي Regular أو Medium حسب رغبتك)
\newfontfamily\boxBodyFont[Script=Arabic]{IBMPlexSansArabic-Light.ttf}

% 2. البيئة المطورة مع لوحة تحكم علوية
\newenvironment{arabicbox}[1]{%
  \begin{center}
    \begin{LTR}
      \begin{tikzpicture}
        \node[
          fill=boxbackground, 
          draw=boxborder, 
          rounded corners=6pt, 
          inner sep=12pt,
          line width=0.6pt % خط إطار أنيق ومتناسق
        ] \bgroup
          \begin{RTL}
            \begin{varwidth}{0.8\linewidth}
              % 1. تطبيق خط العنوان المخصص (بدون الحاجة لـ \textbf الحادة)
              \begin{center}\Large\boxTitleFont{#1}\\[3mm]\end{center} 
              \centering
              \large
              % 2. تطبيق خط المتن المخصص للصندوق
              \boxBodyFont 
}{%
              \vspace{2mm} % مسافة أمان سفلية قبل إغلاق الصندوق
            \end{varwidth}
          \end{RTL}
        \egroup;
      \end{tikzpicture}
    \end{LTR}
  \end{center}
}


\usepackage{titlesec} % حزمة التحكم في تنسيق العناوين

% تخصيص تنسيق القسم (Section) ليصبح بشكل "عرض" (display)
\titleformat{\section}[display] 
    {\fontsize{35}{50}\selectfont\bfseries\filcenter} % حجم العنوان ثابت كما هو
    {\vspace*{50pt}\fontsize{140}{130}\selectfont\thesection} % الحيلة هنا: أضفنا \vspace*{80pt} المحمية فوق الرقم مباشرة
    {0pt}                                             % مسافة صفرية آمنة تمنع تسريب الكلمات
    {} 
    
% --- تنسيق المسافات المحيطة بالقسم ---
\titlespacing*{\section}
    {0pt}    % لا يوجد إزاحة جانبية
    {0pt}    % جعلناها صفر هنا لأننا نقلنا التحكم في المسافة العلوية للأعلى بشكل إجباري ومضمون
    {100pt}  % المسافة الكبيرة تحت العنوان لتفصله عن النص

\usepackage{hyperref} % For hyperlinks in the PDF
% 1. يجب استدعاء حزمة الألوان قبل hyperref للحصول على أسماء ألوان احترافية
\usepackage[dvipsnames]{xcolor} 

% 2. إعدادات الـ hyperref المتقدمة والمحسنة للمشاريع الأكاديمية
\hypersetup{
    unicode=true,          % (حاسم!) يضمن قراءة الحروف العربية في خصائص الـ PDF بدون تشوه
    colorlinks=true,       % تلوين النص بدلاً من المربعات
    % --- تخصيص ألوان أكاديمية رصينة ومريحة للعين ---
    linkcolor=MidnightBlue, % لون الروابط الداخلية (الفهرس، المراجع، المعادلات) - أزرق داكن فخم
    filecolor=Maroon,       % لون روابط الملفات المحلية - نبيذي داكن
    urlcolor=Teal,          % لون الروابط الخارجية للمواقع - أزرق مخضر هادئ
    % --- بيانات الميتا المحقونة في الملف ---
    pdftitle={تحسين جودة القدرة الكهربائية لشبكة كهربائية تغذي حمل غير خطي},
    pdfauthor={محمد علي قطّان},
    pdfsubject={مشروع لمادة التحكم بنظم القدرة}, % إضافة اختيارية متميزة لتحديد نوع المستند
    pdfkeywords={جودة القدرة، مرشحات، تحليل هارمونيك، MATLAB},
    % --- تحسينات تجربة المستخدم عند فتح ملف الـ PDF ---
    pdfstartview={FitH},    % يفتح الملف تلقائياً بحيث يتناسب عرضه مع عرض شاشة القارئ
    bookmarksnumbered=true, % يظهر أرقام الأقسام (1, 2, 3...) في شجرة التنقل الجانبية للـ PDF
    pdfpagemode=UseOutlines % يفتح قائمة التنقل الجانبية (Bookmarks) تلقائياً للمشرف بمجرد فتحه للملف
}